% --- Archon physics formalization source begin ---
% archon:physics
% archon:covers PhyXMiniProblems/problem_phyx_mini_0953.lean
% archon:source-report reports/phyx_mini/problem_phyx_mini_0953.source.json

\chapter{Physics problem phyx\_mini\_0953}
\label{ch:PhyXMiniProblems_problem_phyx_mini_0953}

\paragraph{Problem source.}
\#\# Physical scenario

The lower end of the thin uniform rod is attached to the floor by a frictionless hinge at point \textbackslash{}( P \textbackslash{}). The rod has mass \textbackslash{}( 0.0840\textbackslash{},\textbackslash{}text\{kg\} \textbackslash{}) and length \textbackslash{}( 18.0\textbackslash{},\textbackslash{}text\{cm\} \textbackslash{}) and is in a uniform magnetic field \textbackslash{}( B = 0.120\textbackslash{},\textbackslash{}text\{T\} \textbackslash{}) that is directed into the page. The rod is held at an angle \textbackslash{}( 	heta = 53.0\textasciicircum{}\textbackslash{}circ \textbackslash{}) above the horizontal by a horizontal string that connects the top of the rod to the wall. The rod carries a current \textbackslash{}( I = 12.0\textbackslash{},\textbackslash{}text\{A\} \textbackslash{}) in the direction toward \textbackslash{}( P \textbackslash{}).

\#\# Question

Calculate the tension in the string.

\#\# Answer choices

- A: A: 0.394N

- B: B: 0.472N

- C: C: 0.528N

- D: D: 0.708N

\#\# Figure caption (auxiliary; use the image as primary evidence)

The image depicts a graphical representation commonly used in physics to illustrate the interaction between a current-carrying wire and a magnetic field. Here's a detailed description:

- **Background**: The field is represented by several blue Xs arranged in a grid pattern, indicating a magnetic field directed into the plane of the image.

- **Objects**:
  - A straight, slanted wire carries a current and is illustrated by a tan-colored rod overlaying the field.
  - The current direction in the wire is indicated by a purple arrow labeled "I," pointing downwards along the wire.

- **Angles and Geometry**:
  - The wire is connected at a point labeled "P" on a vertical line.
  - An angle θ (theta) is marked between the vertical line at point P and the direction of the wire, showing the wire’s inclination.

- **Text and Labels**:
  - The magnetic field is denoted with \textbackslash{}(\textbackslash{}vec\{B\}\textbackslash{}) in blue, positioned near the top right.
  - \textbackslash{}(\textbackslash{}theta\textbackslash{}) is marked to indicate the angle between the current I and the vertical reference line.

This image likely represents a scenario involving the magnetic force on a current-carrying conductor in a uniform magnetic field, as described by the right-hand rule or other related principles in electromagnetism.

\#\# Dataset metadata

Magnetism; Physical Model Grounding Reasoning, Multi-Formula Reasoning

\paragraph{Recorded answer/context.}
B: B: 0.472N

\paragraph{Figure/image path.}
/root/proposal\_for\_physic/science-mango/phyx\_mini\_run/phyx\_data/test\_image/953.png

\paragraph{Formalization target.}
create a compiling Lean file with sorry bodies at `PhyXMiniProblems/problem\_phyx\_mini\_0953.lean`. The Lean declarations must preserve the physical quantities, dimensions or dimensional roles, figure labels, governing-law hypotheses, and final relation expressed by this problem.
Use Mathlib/Physlib names found through LeanExplore where available. If a physics API is missing, introduce faithful local abstractions rather than scalar placeholder aliases.

\begin{theorem}[Physics formalization target]
\label{thm:physics:phyx_mini_0953:target}
\lean{PhyXMiniProblems.ProblemPhyXMini0953.string_tension_is_answer_B}
\uses{def:physics:phyx-mini-0953:phyxminiproblems-problemphyxmini0953-forceinnewtons, def:physics:phyx-mini-0953:phyxminiproblems-problemphyxmini0953-hingedcurrentrodsetup, def:physics:phyx-mini-0953:phyxminiproblems-problemphyxmini0953-matcheswrittenhingedrodscenario, def:physics:phyx-mini-0953:phyxminiproblems-problemphyxmini0953-matchesprimaryhingedrodfigure, def:physics:phyx-mini-0953:phyxminiproblems-problemphyxmini0953-matcheshingedrodproblemdata, def:physics:phyx-mini-0953:phyxminiproblems-problemphyxmini0953-usesschoolgravitycalibration, def:physics:phyx-mini-0953:phyxminiproblems-problemphyxmini0953-hasphysicalhingedrodparameters, def:physics:phyx-mini-0953:phyxminiproblems-problemphyxmini0953-hashingedrodgeometry, def:physics:phyx-mini-0953:phyxminiproblems-problemphyxmini0953-hasuniformappliedmagneticfield, def:physics:phyx-mini-0953:phyxminiproblems-problemphyxmini0953-satisfiesmagneticforcelaw, def:physics:phyx-mini-0953:phyxminiproblems-problemphyxmini0953-satisfieshingedrodstaticlaws, def:physics:phyx-mini-0953:phyxminiproblems-problemphyxmini0953-answerchoice, def:physics:phyx-mini-0953:phyxminiproblems-problemphyxmini0953-recordeddatasetanswer, def:physics:phyx-mini-0953:phyxminiproblems-problemphyxmini0953-roundstonearestmillinewton, def:physics:phyx-mini-0953:phyxminiproblems-problemphyxmini0953-isuniquenearestdisplayedtension}
The assigned autoformalize agent should translate this physics problem into Lean declarations in the covered file, with theorem and lemma proof bodies written as `by sorry`.
\end{theorem}
\begin{proof}
This is an autoformalization task, not a proof task. Produce statements that can later be proved by the physics prover without weakening the physical model.
\end{proof}

% --- Archon declaration topology begin ---
\section{Declaration topology}
\begin{definition}[electric Current Dimension]
  \label{def:physics:phyx-mini-0953:phyxminiproblems-problemphyxmini0953-electriccurrentdimension}
  \lean{PhyXMiniProblems.ProblemPhyXMini0953.electricCurrentDimension}
  Electric current has dimension C T⁻¹
\end{definition}
\begin{proof}
  This is the declared model definition of electric Current Dimension.
\end{proof}
\begin{definition}[magnetic Flux Density Dimension]
  \label{def:physics:phyx-mini-0953:phyxminiproblems-problemphyxmini0953-magneticfluxdensitydimension}
  \lean{PhyXMiniProblems.ProblemPhyXMini0953.magneticFluxDensityDimension}
  Magnetic flux density (tesla) has dimension M T⁻¹ C⁻¹
\end{definition}
\begin{proof}
  This is the declared model definition of magnetic Flux Density Dimension.
\end{proof}
\begin{definition}[acceleration Dimension]
  \label{def:physics:phyx-mini-0953:phyxminiproblems-problemphyxmini0953-accelerationdimension}
  \lean{PhyXMiniProblems.ProblemPhyXMini0953.accelerationDimension}
  Acceleration has dimension L T⁻²
\end{definition}
\begin{proof}
  This is the declared model definition of acceleration Dimension.
\end{proof}
\begin{definition}[force Dimension]
  \label{def:physics:phyx-mini-0953:phyxminiproblems-problemphyxmini0953-forcedimension}
  \lean{PhyXMiniProblems.ProblemPhyXMini0953.forceDimension}
  Force (newton) has dimension M L T⁻²
\end{definition}
\begin{proof}
  This is the declared model definition of force Dimension.
\end{proof}
\begin{definition}[torque Dimension]
  \label{def:physics:phyx-mini-0953:phyxminiproblems-problemphyxmini0953-torquedimension}
  \lean{PhyXMiniProblems.ProblemPhyXMini0953.torqueDimension}
  Torque (newton metre) has dimension M L² T⁻²
\end{definition}
\begin{proof}
  This is the declared model definition of torque Dimension.
\end{proof}
\begin{definition}[Mass Quantity]
  \label{def:physics:phyx-mini-0953:phyxminiproblems-problemphyxmini0953-massquantity}
  \lean{PhyXMiniProblems.ProblemPhyXMini0953.MassQuantity}
  A nonnegative, unit-independent physical mass
\end{definition}
\begin{proof}
  This is the declared model definition of Mass Quantity.
\end{proof}
\begin{definition}[Length Quantity]
  \label{def:physics:phyx-mini-0953:phyxminiproblems-problemphyxmini0953-lengthquantity}
  \lean{PhyXMiniProblems.ProblemPhyXMini0953.LengthQuantity}
  A nonnegative, unit-independent physical length
\end{definition}
\begin{proof}
  This is the declared model definition of Length Quantity.
\end{proof}
\begin{definition}[Electric Current Magnitude]
  \label{def:physics:phyx-mini-0953:phyxminiproblems-problemphyxmini0953-electriccurrentmagnitude}
  \lean{PhyXMiniProblems.ProblemPhyXMini0953.ElectricCurrentMagnitude}
  \uses{def:physics:phyx-mini-0953:phyxminiproblems-problemphyxmini0953-electriccurrentdimension}
  A nonnegative, unit-independent current magnitude
\end{definition}
\begin{proof}
  This is the declared model definition of Electric Current Magnitude.
\end{proof}
\begin{definition}[Magnetic Flux Density Magnitude]
  \label{def:physics:phyx-mini-0953:phyxminiproblems-problemphyxmini0953-magneticfluxdensitymagnitude}
  \lean{PhyXMiniProblems.ProblemPhyXMini0953.MagneticFluxDensityMagnitude}
  \uses{def:physics:phyx-mini-0953:phyxminiproblems-problemphyxmini0953-magneticfluxdensitydimension}
  A nonnegative, unit-independent magnetic-flux-density magnitude
\end{definition}
\begin{proof}
  This is the declared model definition of Magnetic Flux Density Magnitude.
\end{proof}
\begin{definition}[Acceleration Magnitude]
  \label{def:physics:phyx-mini-0953:phyxminiproblems-problemphyxmini0953-accelerationmagnitude}
  \lean{PhyXMiniProblems.ProblemPhyXMini0953.AccelerationMagnitude}
  \uses{def:physics:phyx-mini-0953:phyxminiproblems-problemphyxmini0953-accelerationdimension}
  A nonnegative, unit-independent acceleration magnitude
\end{definition}
\begin{proof}
  This is the declared model definition of Acceleration Magnitude.
\end{proof}
\begin{definition}[Force Magnitude]
  \label{def:physics:phyx-mini-0953:phyxminiproblems-problemphyxmini0953-forcemagnitude}
  \lean{PhyXMiniProblems.ProblemPhyXMini0953.ForceMagnitude}
  \uses{def:physics:phyx-mini-0953:phyxminiproblems-problemphyxmini0953-forcedimension}
  A nonnegative, unit-independent force magnitude
\end{definition}
\begin{proof}
  This is the declared model definition of Force Magnitude.
\end{proof}
\begin{definition}[Torque Magnitude]
  \label{def:physics:phyx-mini-0953:phyxminiproblems-problemphyxmini0953-torquemagnitude}
  \lean{PhyXMiniProblems.ProblemPhyXMini0953.TorqueMagnitude}
  \uses{def:physics:phyx-mini-0953:phyxminiproblems-problemphyxmini0953-torquedimension}
  A nonnegative, unit-independent moment magnitude about the hinge
\end{definition}
\begin{proof}
  This is the declared model definition of Torque Magnitude.
\end{proof}
\begin{definition}[nonnegative SIReadout]
  \label{def:physics:phyx-mini-0953:phyxminiproblems-problemphyxmini0953-nonnegativesireadout}
  \lean{PhyXMiniProblems.ProblemPhyXMini0953.nonnegativeSIReadout}
  Read a nonnegative dimensionful quantity in coherent SI base units
\end{definition}
\begin{proof}
  This is the declared model definition of nonnegative SIReadout.
\end{proof}
\begin{definition}[mass In Kilograms]
  \label{def:physics:phyx-mini-0953:phyxminiproblems-problemphyxmini0953-massinkilograms}
  \lean{PhyXMiniProblems.ProblemPhyXMini0953.massInKilograms}
  \uses{def:physics:phyx-mini-0953:phyxminiproblems-problemphyxmini0953-massquantity, def:physics:phyx-mini-0953:phyxminiproblems-problemphyxmini0953-nonnegativesireadout}
  Read a mass in kilograms
\end{definition}
\begin{proof}
  This is the declared model definition of mass In Kilograms.
\end{proof}
\begin{definition}[length In Meters]
  \label{def:physics:phyx-mini-0953:phyxminiproblems-problemphyxmini0953-lengthinmeters}
  \lean{PhyXMiniProblems.ProblemPhyXMini0953.lengthInMeters}
  \uses{def:physics:phyx-mini-0953:phyxminiproblems-problemphyxmini0953-lengthquantity, def:physics:phyx-mini-0953:phyxminiproblems-problemphyxmini0953-nonnegativesireadout}
  Read a length in metres
\end{definition}
\begin{proof}
  This is the declared model definition of length In Meters.
\end{proof}
\begin{definition}[length In Centimeters]
  \label{def:physics:phyx-mini-0953:phyxminiproblems-problemphyxmini0953-lengthincentimeters}
  \lean{PhyXMiniProblems.ProblemPhyXMini0953.lengthInCentimeters}
  \uses{def:physics:phyx-mini-0953:phyxminiproblems-problemphyxmini0953-lengthquantity, def:physics:phyx-mini-0953:phyxminiproblems-problemphyxmini0953-lengthinmeters}
  Read a length in centimetres
\end{definition}
\begin{proof}
  This is the declared model definition of length In Centimeters.
\end{proof}
\begin{definition}[current In Amperes]
  \label{def:physics:phyx-mini-0953:phyxminiproblems-problemphyxmini0953-currentinamperes}
  \lean{PhyXMiniProblems.ProblemPhyXMini0953.currentInAmperes}
  \uses{def:physics:phyx-mini-0953:phyxminiproblems-problemphyxmini0953-electriccurrentmagnitude, def:physics:phyx-mini-0953:phyxminiproblems-problemphyxmini0953-nonnegativesireadout}
  Read a current magnitude in amperes
\end{definition}
\begin{proof}
  This is the declared model definition of current In Amperes.
\end{proof}
\begin{definition}[magnetic Flux Density In Teslas]
  \label{def:physics:phyx-mini-0953:phyxminiproblems-problemphyxmini0953-magneticfluxdensityinteslas}
  \lean{PhyXMiniProblems.ProblemPhyXMini0953.magneticFluxDensityInTeslas}
  \uses{def:physics:phyx-mini-0953:phyxminiproblems-problemphyxmini0953-magneticfluxdensitymagnitude, def:physics:phyx-mini-0953:phyxminiproblems-problemphyxmini0953-nonnegativesireadout}
  Read a magnetic flux density in teslas
\end{definition}
\begin{proof}
  This is the declared model definition of magnetic Flux Density In Teslas.
\end{proof}
\begin{definition}[acceleration In Meters Per Second Squared]
  \label{def:physics:phyx-mini-0953:phyxminiproblems-problemphyxmini0953-accelerationinmeterspersecondsquared}
  \lean{PhyXMiniProblems.ProblemPhyXMini0953.accelerationInMetersPerSecondSquared}
  \uses{def:physics:phyx-mini-0953:phyxminiproblems-problemphyxmini0953-accelerationmagnitude, def:physics:phyx-mini-0953:phyxminiproblems-problemphyxmini0953-nonnegativesireadout}
  Read an acceleration magnitude in metres per second squared
\end{definition}
\begin{proof}
  This is the declared model definition of acceleration In Meters Per Second Squared.
\end{proof}
\begin{definition}[force In Newtons]
  \label{def:physics:phyx-mini-0953:phyxminiproblems-problemphyxmini0953-forceinnewtons}
  \lean{PhyXMiniProblems.ProblemPhyXMini0953.forceInNewtons}
  \uses{def:physics:phyx-mini-0953:phyxminiproblems-problemphyxmini0953-forcemagnitude, def:physics:phyx-mini-0953:phyxminiproblems-problemphyxmini0953-nonnegativesireadout}
  Read a force magnitude in newtons
\end{definition}
\begin{proof}
  This is the declared model definition of force In Newtons.
\end{proof}
\begin{definition}[torque In Newton Meters]
  \label{def:physics:phyx-mini-0953:phyxminiproblems-problemphyxmini0953-torqueinnewtonmeters}
  \lean{PhyXMiniProblems.ProblemPhyXMini0953.torqueInNewtonMeters}
  \uses{def:physics:phyx-mini-0953:phyxminiproblems-problemphyxmini0953-torquemagnitude, def:physics:phyx-mini-0953:phyxminiproblems-problemphyxmini0953-nonnegativesireadout}
  Read a moment magnitude in newton metres
\end{definition}
\begin{proof}
  This is the declared model definition of torque In Newton Meters.
\end{proof}
\begin{definition}[degrees To Angle]
  \label{def:physics:phyx-mini-0953:phyxminiproblems-problemphyxmini0953-degreestoangle}
  \lean{PhyXMiniProblems.ProblemPhyXMini0953.degreesToAngle}
  Convert a degree readout to Mathlib's quotient type of real angles
\end{definition}
\begin{proof}
  This is the declared model definition of degrees To Angle.
\end{proof}
\begin{definition}[Spatial Vector]
  \label{def:physics:phyx-mini-0953:phyxminiproblems-problemphyxmini0953-spatialvector}
  \lean{PhyXMiniProblems.ProblemPhyXMini0953.SpatialVector}
  A coherent-SI direction vector in physical three-space
\end{definition}
\begin{proof}
  This is the declared model definition of Spatial Vector.
\end{proof}
\begin{definition}[rightward Unit Vector]
  \label{def:physics:phyx-mini-0953:phyxminiproblems-problemphyxmini0953-rightwardunitvector}
  \lean{PhyXMiniProblems.ProblemPhyXMini0953.rightwardUnitVector}
  \uses{def:physics:phyx-mini-0953:phyxminiproblems-problemphyxmini0953-spatialvector}
  Rightward unit vector in the plane of the supplied image
\end{definition}
\begin{proof}
  This is the declared model definition of rightward Unit Vector.
\end{proof}
\begin{definition}[upward Unit Vector]
  \label{def:physics:phyx-mini-0953:phyxminiproblems-problemphyxmini0953-upwardunitvector}
  \lean{PhyXMiniProblems.ProblemPhyXMini0953.upwardUnitVector}
  \uses{def:physics:phyx-mini-0953:phyxminiproblems-problemphyxmini0953-spatialvector}
  Upward unit vector in the plane of the supplied image
\end{definition}
\begin{proof}
  This is the declared model definition of upward Unit Vector.
\end{proof}
\begin{definition}[into Page Unit Vector]
  \label{def:physics:phyx-mini-0953:phyxminiproblems-problemphyxmini0953-intopageunitvector}
  \lean{PhyXMiniProblems.ProblemPhyXMini0953.intoPageUnitVector}
  \uses{def:physics:phyx-mini-0953:phyxminiproblems-problemphyxmini0953-spatialvector}
  Unit vector into the page, represented by the field-cross convention
\end{definition}
\begin{proof}
  This is the declared model definition of into Page Unit Vector.
\end{proof}
\begin{definition}[spatial Cross]
  \label{def:physics:phyx-mini-0953:phyxminiproblems-problemphyxmini0953-spatialcross}
  \lean{PhyXMiniProblems.ProblemPhyXMini0953.spatialCross}
  \uses{def:physics:phyx-mini-0953:phyxminiproblems-problemphyxmini0953-spatialvector}
  The ordinary right-handed cross product on Euclidean three-vectors
\end{definition}
\begin{proof}
  This is the declared model definition of spatial Cross.
\end{proof}
\begin{definition}[rod Direction Above Horizontal]
  \label{def:physics:phyx-mini-0953:phyxminiproblems-problemphyxmini0953-roddirectionabovehorizontal}
  \lean{PhyXMiniProblems.ProblemPhyXMini0953.rodDirectionAboveHorizontal}
  \uses{def:physics:phyx-mini-0953:phyxminiproblems-problemphyxmini0953-spatialvector, def:physics:phyx-mini-0953:phyxminiproblems-problemphyxmini0953-rightwardunitvector, def:physics:phyx-mini-0953:phyxminiproblems-problemphyxmini0953-upwardunitvector}
  Unit direction of a rod inclined by angle above the horizontal
\end{definition}
\begin{proof}
  This is the declared model definition of rod Direction Above Horizontal.
\end{proof}
\begin{definition}[Rod Mass Model]
  \label{def:physics:phyx-mini-0953:phyxminiproblems-problemphyxmini0953-rodmassmodel}
  \lean{PhyXMiniProblems.ProblemPhyXMini0953.RodMassModel}
  Idealized distribution of mass along the rod
\end{definition}
\begin{proof}
  This is the declared model definition of Rod Mass Model.
\end{proof}
\begin{definition}[Hinge Model]
  \label{def:physics:phyx-mini-0953:phyxminiproblems-problemphyxmini0953-hingemodel}
  \lean{PhyXMiniProblems.ProblemPhyXMini0953.HingeModel}
  Mechanical model of the attachment at the lower endpoint
\end{definition}
\begin{proof}
  This is the declared model definition of Hinge Model.
\end{proof}
\begin{definition}[String Orientation]
  \label{def:physics:phyx-mini-0953:phyxminiproblems-problemphyxmini0953-stringorientation}
  \lean{PhyXMiniProblems.ProblemPhyXMini0953.StringOrientation}
  Orientation class of the string shown at the rod's upper endpoint
\end{definition}
\begin{proof}
  This is the declared model definition of String Orientation.
\end{proof}
\begin{definition}[String Attachment]
  \label{def:physics:phyx-mini-0953:phyxminiproblems-problemphyxmini0953-stringattachment}
  \lean{PhyXMiniProblems.ProblemPhyXMini0953.StringAttachment}
  Point of the rod at which the supporting string is attached
\end{definition}
\begin{proof}
  This is the declared model definition of String Attachment.
\end{proof}
\begin{definition}[Hinge Support Surface]
  \label{def:physics:phyx-mini-0953:phyxminiproblems-problemphyxmini0953-hingesupportsurface}
  \lean{PhyXMiniProblems.ProblemPhyXMini0953.HingeSupportSurface}
  Surface carrying the hinge at P
\end{definition}
\begin{proof}
  This is the declared model definition of Hinge Support Surface.
\end{proof}
\begin{definition}[Diagram Direction]
  \label{def:physics:phyx-mini-0953:phyxminiproblems-problemphyxmini0953-diagramdirection}
  \lean{PhyXMiniProblems.ProblemPhyXMini0953.DiagramDirection}
  Coarse directions sufficient to transcribe the supplied raster
\end{definition}
\begin{proof}
  This is the declared model definition of Diagram Direction.
\end{proof}
\begin{definition}[Hinged Rod Figure]
  \label{def:physics:phyx-mini-0953:phyxminiproblems-problemphyxmini0953-hingedrodfigure}
  \lean{PhyXMiniProblems.ProblemPhyXMini0953.HingedRodFigure}
  \uses{def:physics:phyx-mini-0953:phyxminiproblems-problemphyxmini0953-diagramdirection}
  Literal labels and qualitative geometry visible in image 953.png
\end{definition}
\begin{proof}
  This is the declared model definition of Hinged Rod Figure.
\end{proof}
\begin{definition}[Hinged Current Rod Setup]
  \label{def:physics:phyx-mini-0953:phyxminiproblems-problemphyxmini0953-hingedcurrentrodsetup}
  \lean{PhyXMiniProblems.ProblemPhyXMini0953.HingedCurrentRodSetup}
  \uses{def:physics:phyx-mini-0953:phyxminiproblems-problemphyxmini0953-massquantity, def:physics:phyx-mini-0953:phyxminiproblems-problemphyxmini0953-lengthquantity, def:physics:phyx-mini-0953:phyxminiproblems-problemphyxmini0953-electriccurrentmagnitude, def:physics:phyx-mini-0953:phyxminiproblems-problemphyxmini0953-magneticfluxdensitymagnitude, def:physics:phyx-mini-0953:phyxminiproblems-problemphyxmini0953-accelerationmagnitude, def:physics:phyx-mini-0953:phyxminiproblems-problemphyxmini0953-forcemagnitude, def:physics:phyx-mini-0953:phyxminiproblems-problemphyxmini0953-torquemagnitude, def:physics:phyx-mini-0953:phyxminiproblems-problemphyxmini0953-spatialvector, def:physics:phyx-mini-0953:phyxminiproblems-problemphyxmini0953-rodmassmodel, def:physics:phyx-mini-0953:phyxminiproblems-problemphyxmini0953-hingemodel, def:physics:phyx-mini-0953:phyxminiproblems-problemphyxmini0953-stringorientation, def:physics:phyx-mini-0953:phyxminiproblems-problemphyxmini0953-stringattachment, def:physics:phyx-mini-0953:phyxminiproblems-problemphyxmini0953-hingesupportsurface, def:physics:phyx-mini-0953:phyxminiproblems-problemphyxmini0953-hingedrodfigure}
  Independent observables for the supported rod. In particular, the tension and the four moment magnitudes are stored independently and are constrained only by the governing laws below
\end{definition}
\begin{proof}
  This is the declared model definition of Hinged Current Rod Setup.
\end{proof}
\begin{definition}[Matches Written Hinged Rod Scenario]
  \label{def:physics:phyx-mini-0953:phyxminiproblems-problemphyxmini0953-matcheswrittenhingedrodscenario}
  \lean{PhyXMiniProblems.ProblemPhyXMini0953.MatchesWrittenHingedRodScenario}
  \uses{def:physics:phyx-mini-0953:phyxminiproblems-problemphyxmini0953-hingedcurrentrodsetup}
  The idealizations stated in the written physical scenario
\end{definition}
\begin{proof}
  This is the declared model definition of Matches Written Hinged Rod Scenario.
\end{proof}
\begin{definition}[Matches Primary Hinged Rod Figure]
  \label{def:physics:phyx-mini-0953:phyxminiproblems-problemphyxmini0953-matchesprimaryhingedrodfigure}
  \lean{PhyXMiniProblems.ProblemPhyXMini0953.MatchesPrimaryHingedRodFigure}
  \uses{def:physics:phyx-mini-0953:phyxminiproblems-problemphyxmini0953-hingedcurrentrodsetup}
  Primary-raster evidence. The image shows the angle arc between the rod and the horizontal baseline at P, a downward current arrow along the rod, field crosses, and a horizontal string extending left from the upper endpoint
\end{definition}
\begin{proof}
  This is the declared model definition of Matches Primary Hinged Rod Figure.
\end{proof}
\begin{definition}[Matches Hinged Rod Problem Data]
  \label{def:physics:phyx-mini-0953:phyxminiproblems-problemphyxmini0953-matcheshingedrodproblemdata}
  \lean{PhyXMiniProblems.ProblemPhyXMini0953.MatchesHingedRodProblemData}
  \uses{def:physics:phyx-mini-0953:phyxminiproblems-problemphyxmini0953-massinkilograms, def:physics:phyx-mini-0953:phyxminiproblems-problemphyxmini0953-lengthincentimeters, def:physics:phyx-mini-0953:phyxminiproblems-problemphyxmini0953-currentinamperes, def:physics:phyx-mini-0953:phyxminiproblems-problemphyxmini0953-magneticfluxdensityinteslas, def:physics:phyx-mini-0953:phyxminiproblems-problemphyxmini0953-degreestoangle, def:physics:phyx-mini-0953:phyxminiproblems-problemphyxmini0953-hingedcurrentrodsetup}
  Exact stated measurements, expressed in their named units
\end{definition}
\begin{proof}
  This is the declared model definition of Matches Hinged Rod Problem Data.
\end{proof}
\begin{definition}[Uses School Gravity Calibration]
  \label{def:physics:phyx-mini-0953:phyxminiproblems-problemphyxmini0953-usesschoolgravitycalibration}
  \lean{PhyXMiniProblems.ProblemPhyXMini0953.UsesSchoolGravityCalibration}
  \uses{def:physics:phyx-mini-0953:phyxminiproblems-problemphyxmini0953-accelerationinmeterspersecondsquared, def:physics:phyx-mini-0953:phyxminiproblems-problemphyxmini0953-hingedcurrentrodsetup}
  Rounded gravitational acceleration convention used by the answer choices
\end{definition}
\begin{proof}
  This is the declared model definition of Uses School Gravity Calibration.
\end{proof}
\begin{definition}[Has Physical Hinged Rod Parameters]
  \label{def:physics:phyx-mini-0953:phyxminiproblems-problemphyxmini0953-hasphysicalhingedrodparameters}
  \lean{PhyXMiniProblems.ProblemPhyXMini0953.HasPhysicalHingedRodParameters}
  \uses{def:physics:phyx-mini-0953:phyxminiproblems-problemphyxmini0953-massinkilograms, def:physics:phyx-mini-0953:phyxminiproblems-problemphyxmini0953-lengthinmeters, def:physics:phyx-mini-0953:phyxminiproblems-problemphyxmini0953-currentinamperes, def:physics:phyx-mini-0953:phyxminiproblems-problemphyxmini0953-magneticfluxdensityinteslas, def:physics:phyx-mini-0953:phyxminiproblems-problemphyxmini0953-accelerationinmeterspersecondsquared, def:physics:phyx-mini-0953:phyxminiproblems-problemphyxmini0953-forceinnewtons, def:physics:phyx-mini-0953:phyxminiproblems-problemphyxmini0953-hingedcurrentrodsetup}
  Positivity and nondegeneracy conditions for the physical branch
\end{definition}
\begin{proof}
  This is the declared model definition of Has Physical Hinged Rod Parameters.
\end{proof}
\begin{definition}[Has Hinged Rod Geometry]
  \label{def:physics:phyx-mini-0953:phyxminiproblems-problemphyxmini0953-hashingedrodgeometry}
  \lean{PhyXMiniProblems.ProblemPhyXMini0953.HasHingedRodGeometry}
  \uses{def:physics:phyx-mini-0953:phyxminiproblems-problemphyxmini0953-rightwardunitvector, def:physics:phyx-mini-0953:phyxminiproblems-problemphyxmini0953-upwardunitvector, def:physics:phyx-mini-0953:phyxminiproblems-problemphyxmini0953-intopageunitvector, def:physics:phyx-mini-0953:phyxminiproblems-problemphyxmini0953-roddirectionabovehorizontal, def:physics:phyx-mini-0953:phyxminiproblems-problemphyxmini0953-hingedcurrentrodsetup}
  The physical direction vectors encode the geometry described by the prose and confirmed by the raster. Current is opposite the rod axis because it flows toward P; tension points toward the wall
\end{definition}
\begin{proof}
  This is the declared model definition of Has Hinged Rod Geometry.
\end{proof}
\begin{definition}[Has Uniform Applied Magnetic Field]
  \label{def:physics:phyx-mini-0953:phyxminiproblems-problemphyxmini0953-hasuniformappliedmagneticfield}
  \lean{PhyXMiniProblems.ProblemPhyXMini0953.HasUniformAppliedMagneticField}
  \uses{def:physics:phyx-mini-0953:phyxminiproblems-problemphyxmini0953-magneticfluxdensityinteslas, def:physics:phyx-mini-0953:phyxminiproblems-problemphyxmini0953-hingedcurrentrodsetup}
  The Physlib magnetic field has the stated direction and coherent-SI magnitude throughout the modeled uniform-field region at the observation time
\end{definition}
\begin{proof}
  This is the declared model definition of Has Uniform Applied Magnetic Field.
\end{proof}
\begin{definition}[Satisfies Magnetic Force Law]
  \label{def:physics:phyx-mini-0953:phyxminiproblems-problemphyxmini0953-satisfiesmagneticforcelaw}
  \lean{PhyXMiniProblems.ProblemPhyXMini0953.SatisfiesMagneticForceLaw}
  \uses{def:physics:phyx-mini-0953:phyxminiproblems-problemphyxmini0953-lengthinmeters, def:physics:phyx-mini-0953:phyxminiproblems-problemphyxmini0953-currentinamperes, def:physics:phyx-mini-0953:phyxminiproblems-problemphyxmini0953-magneticfluxdensityinteslas, def:physics:phyx-mini-0953:phyxminiproblems-problemphyxmini0953-forceinnewtons, def:physics:phyx-mini-0953:phyxminiproblems-problemphyxmini0953-spatialcross, def:physics:phyx-mini-0953:phyxminiproblems-problemphyxmini0953-hingedcurrentrodsetup}
  For a straight rod perpendicular to a uniform field, the magnetic resultant has magnitude I L B; its direction is the right-handed cross product of the current direction with the field direction
\end{definition}
\begin{proof}
  This is the declared model definition of Satisfies Magnetic Force Law.
\end{proof}
\begin{definition}[Satisfies Hinged Rod Static Laws]
  \label{def:physics:phyx-mini-0953:phyxminiproblems-problemphyxmini0953-satisfieshingedrodstaticlaws}
  \lean{PhyXMiniProblems.ProblemPhyXMini0953.SatisfiesHingedRodStaticLaws}
  \uses{def:physics:phyx-mini-0953:phyxminiproblems-problemphyxmini0953-massinkilograms, def:physics:phyx-mini-0953:phyxminiproblems-problemphyxmini0953-lengthinmeters, def:physics:phyx-mini-0953:phyxminiproblems-problemphyxmini0953-accelerationinmeterspersecondsquared, def:physics:phyx-mini-0953:phyxminiproblems-problemphyxmini0953-forceinnewtons, def:physics:phyx-mini-0953:phyxminiproblems-problemphyxmini0953-torqueinnewtonmeters, def:physics:phyx-mini-0953:phyxminiproblems-problemphyxmini0953-hingedcurrentrodsetup}
  Weight, geometric moment arms, the frictionless-hinge condition, and static rotational equilibrium about P. Uniform gravity and magnetic loading place their resultants at the rod midpoint. These are general balance laws and do not state the requested tension formula or numerical answer
\end{definition}
\begin{proof}
  This is the declared model definition of Satisfies Hinged Rod Static Laws.
\end{proof}
\begin{lemma}[tension formula]
  \label{lem:physics:phyx-mini-0953:phyxminiproblems-problemphyxmini0953-tension-formula}
  \lean{PhyXMiniProblems.ProblemPhyXMini0953.tension_formula}
  \uses{def:physics:phyx-mini-0953:phyxminiproblems-problemphyxmini0953-massinkilograms, def:physics:phyx-mini-0953:phyxminiproblems-problemphyxmini0953-lengthinmeters, def:physics:phyx-mini-0953:phyxminiproblems-problemphyxmini0953-currentinamperes, def:physics:phyx-mini-0953:phyxminiproblems-problemphyxmini0953-magneticfluxdensityinteslas, def:physics:phyx-mini-0953:phyxminiproblems-problemphyxmini0953-accelerationinmeterspersecondsquared, def:physics:phyx-mini-0953:phyxminiproblems-problemphyxmini0953-forceinnewtons, def:physics:phyx-mini-0953:phyxminiproblems-problemphyxmini0953-hingedcurrentrodsetup, def:physics:phyx-mini-0953:phyxminiproblems-problemphyxmini0953-hasphysicalhingedrodparameters, def:physics:phyx-mini-0953:phyxminiproblems-problemphyxmini0953-satisfiesmagneticforcelaw, def:physics:phyx-mini-0953:phyxminiproblems-problemphyxmini0953-satisfieshingedrodstaticlaws}
  Eliminating the separately modeled forces and moments gives the general tension formula for this geometry. This is a derived conclusion, not a law field
\end{lemma}
\begin{proof}
  The conclusion follows from the governing relations and figure readouts named in the statement of tension formula.
\end{proof}
\begin{definition}[Answer Choice]
  \label{def:physics:phyx-mini-0953:phyxminiproblems-problemphyxmini0953-answerchoice}
  \lean{PhyXMiniProblems.ProblemPhyXMini0953.AnswerChoice}
  The four tension choices printed in the source
\end{definition}
\begin{proof}
  This is the declared model definition of Answer Choice.
\end{proof}
\begin{definition}[tension In Newtons]
  \label{def:physics:phyx-mini-0953:phyxminiproblems-problemphyxmini0953-answerchoice-tensioninnewtons}
  \lean{PhyXMiniProblems.ProblemPhyXMini0953.AnswerChoice.tensionInNewtons}
  \uses{def:physics:phyx-mini-0953:phyxminiproblems-problemphyxmini0953-answerchoice}
  Tension in newtons printed beside each answer label
\end{definition}
\begin{proof}
  This is the declared model definition of tension In Newtons.
\end{proof}
\begin{definition}[recorded Dataset Answer]
  \label{def:physics:phyx-mini-0953:phyxminiproblems-problemphyxmini0953-recordeddatasetanswer}
  \lean{PhyXMiniProblems.ProblemPhyXMini0953.recordedDatasetAnswer}
  \uses{def:physics:phyx-mini-0953:phyxminiproblems-problemphyxmini0953-answerchoice}
  The answer label recorded in the supplied dataset metadata
\end{definition}
\begin{proof}
  This is the declared model definition of recorded Dataset Answer.
\end{proof}
\begin{definition}[Rounds To Nearest Milli Newton]
  \label{def:physics:phyx-mini-0953:phyxminiproblems-problemphyxmini0953-roundstonearestmillinewton}
  \lean{PhyXMiniProblems.ProblemPhyXMini0953.RoundsToNearestMilliNewton}
  The actual tension rounds to a displayed value at the nearest 0.001 N
\end{definition}
\begin{proof}
  This is the declared model definition of Rounds To Nearest Milli Newton.
\end{proof}
\begin{definition}[Is Unique Nearest Displayed Tension]
  \label{def:physics:phyx-mini-0953:phyxminiproblems-problemphyxmini0953-isuniquenearestdisplayedtension}
  \lean{PhyXMiniProblems.ProblemPhyXMini0953.IsUniqueNearestDisplayedTension}
  \uses{def:physics:phyx-mini-0953:phyxminiproblems-problemphyxmini0953-answerchoice}
  A displayed choice is strictly nearer than every alternative choice
\end{definition}
\begin{proof}
  This is the declared model definition of Is Unique Nearest Displayed Tension.
\end{proof}
% --- Archon declaration topology end ---
% --- Archon physics formalization source end ---
